% GENERATED FILE. Edit canonical lessons, then run npm run generate:books.
% canonical-source-hash: fnv1a64:f3b4736254ea83b6
% canonical-lessons: SA-C261-avam, SA-C261-yuvam, SA-C261-dual-object-practice

\chapter{Us two and you two as objects}
\label{ch:sa-us-two-you-two-objects}
\hlchaptermodality{261}

\begin{chapteropening}
\textbf{By the end of this chapter:} \emph{I can distinguish us two from you two when the pair is the object.}

Complete chapter 261 by deciding whether an object pair includes the speaker and producing the matching dual form.
\end{chapteropening}

\section[\texorpdfstring{\sk{आवाम्} (āvām)}{आवाम् (āvām)}]{\sk{आवाम्} (āvām) — us two, as the object}

\label{lesson:SA-C261-avam}



\begin{quote}
Say \textbf{\sk{माम्}}, “me”. Then recall the Sanskrit number reserved for exactly two: the dual. Now make the object be the speaker and one other person.
\end{quote}

\subsection*{You'll want to know: \sk{आवाम्}}
\textbf{\sk{आवाम्}} (\emph{āvām}) means “us two” as the object.

\begin{quote}
\textbf{\sk{सः} \sk{आवाम्} \sk{पश्यति।}}
\end{quote}

\begin{quote}
\emph{saḥ āvām paśyati.}
\end{quote}

\begin{quote}
\textbf{He sees us two.}
\end{quote}

Whitney §491 lists \textbf{\sk{आवाम्}} as the independent dual form used for both subject and object. Here its position in the known frame makes the object job clear. Hold the first vowel long: \emph{ā-vām}.

\begin{cousinweb}
As a grammatical form, \textbf{\sk{आवाम्}} comes from the first-person pronoun paradigm; it is not a word from a separate lexical root. Whitney's \emph{Sanskrit Grammar} §491 records it as the independent dual used for both subject and object. This lesson makes no deeper historical derivation claim.

Source: Whitney, \emph{Sanskrit Grammar} (1889), §491.
\end{cousinweb}

\subsection*{Your turn}
Say these aloud:

\begin{itemize}
\raggedright
  \item one speaker as object — \emph{mām}
  \item the speaker plus one other as object — \emph{āvām}
  \item “he sees us two” — \emph{saḥ āvām paśyati}
\end{itemize}

\begin{tcolorbox}[breakable,colback=teal!4,colframe=teal!35!black,title={Before you move on}]
The speaker and exactly one companion are seen. Which object form means “us two”? (\textbf{\sk{आवाम्}}, \emph{āvām}.)
\end{tcolorbox}
\section[\texorpdfstring{\sk{युवाम्} (yuvām)}{युवाम् (yuvām)}]{\sk{युवाम्} (yuvām) — you two, as the object}

\label{lesson:SA-C261-yuvam}



\begin{quote}
One familiar addressee as object is \textbf{\sk{त्वाम्}} (\emph{tvām}). The speaker plus one other is \textbf{\sk{आवाम्}} (\emph{āvām}). Turn the dual toward two addressees.
\end{quote}

\subsection*{You'll want to know: \sk{युवाम्}}
\textbf{\sk{युवाम्}} (\emph{yuvām}) means “you two” as the object.

\begin{quote}
\textbf{\sk{सः} \sk{युवाम्} \sk{पश्यति।}}
\end{quote}

\begin{quote}
\emph{saḥ yuvām paśyati.}
\end{quote}

\begin{quote}
\textbf{He sees you two.}
\end{quote}

Whitney §491 lists \textbf{\sk{युवाम्}} beside \textbf{\sk{आवाम्}}. Both are independent dual forms; the opening tells you whether the pair includes the speaker or consists of the two people addressed.

\begin{cousinweb}
As a grammatical form, \textbf{\sk{युवाम्}} comes from the second-person pronoun paradigm; it is not a word from a separate lexical root. Whitney's \emph{Sanskrit Grammar} §491 records it as the independent dual used for both subject and object. This lesson makes no deeper historical derivation claim.

Source: Whitney, \emph{Sanskrit Grammar} (1889), §491.
\end{cousinweb}

\subsection*{Your turn}
Say these aloud:

\begin{itemize}
\raggedright
  \item “he sees us two” — \emph{saḥ āvām paśyati}
  \item “he sees you two” — \emph{saḥ yuvām paśyati}
  \item our pair — \emph{āvām}; your pair — \emph{yuvām}
\end{itemize}

\begin{tcolorbox}[breakable,colback=teal!4,colframe=teal!35!black,title={Before you move on}]
Exactly two familiar addressees are seen. Which object form do you use? (\textbf{\sk{युवाम्}}, \emph{yuvām}.)
\end{tcolorbox}
\section[Practice]{Does the pair include the speaker?}

\label{lesson:SA-C261-dual-object-practice}



\begin{quote}
Give “us two”, then “you two”. (\emph{Āvām; yuvām.})
\end{quote}

\begin{grammarlens}[title={Grammar Lens: choose the pair}]
Both forms point to exactly two people. Choose by whether the speaker belongs to the pair:

\begin{itemize}
\raggedright
  \item speaker plus one other: \textbf{\sk{आवाम्}} (\emph{āvām}), “us two”;
  \item two familiar addressees: \textbf{\sk{युवाम्}} (\emph{yuvām}), “you two”.
\end{itemize}

That is the dual contrast. Count two first; then ask whether the speaker is inside the pair.
\end{grammarlens}

\subsection*{Your turn}
\begin{itemize}
\raggedright
  \item \emph{Say it:} “he sees us two” — \emph{saḥ āvām paśyati}
  \item \emph{Say it:} “he sees you two” — \emph{saḥ yuvām paśyati}
  \item \emph{Contrast:} “me” then “us two” — \emph{mām; āvām}
  \item \emph{Contrast:} “you one” then “you two” — \emph{tvām; yuvām}
  \item \emph{Say it:} pair includes the speaker — \emph{āvām}; pair is addressed — \emph{yuvām}
  \item \emph{Contrast:} \emph{āvām; yuvām; āvām; yuvām}
\end{itemize}

\begin{tcolorbox}[breakable,colback=teal!4,colframe=teal!35!black,title={Before you move on}]
Say both dual sentences, then say the deciding question. (\emph{Saḥ āvām paśyati; saḥ yuvām paśyati; does the pair include the speaker?})
\end{tcolorbox}
